\documentclass[12pt,letterpaper]{article}

% --- Paquetes Requeridos ---
\usepackage[spanish]{babel}
\usepackage[utf8]{inputenc}
\usepackage{geometry}
\geometry{margin=1in}
\usepackage{hyperref}
\usepackage{listings}
\usepackage{xcolor}
\usepackage{booktabs}
\usepackage{fancyhdr}
\usepackage{amsmath}
\usepackage{graphicx}

% --- Configuración de Enlaces ---
\hypersetup{
    colorlinks=true,
    linkcolor=blue,
    filecolor=magenta,      
    urlcolor=cyan,
    pdftitle={Proyecto 2 - DLP: Generador YALex/YAPar},
    pdfpagemode=FullScreen,
}

% --- Configuración de listings (Código) ---
\definecolor{codegreen}{rgb}{0,0.6,0}
\definecolor{codegray}{rgb}{0.5,0.5,0.5}
\definecolor{codepurple}{rgb}{0.58,0,0.82}
\definecolor{backcolour}{rgb}{0.95,0.95,0.92}

\lstdefinestyle{mystyle}{
    backgroundcolor=\color{backcolour},   
    commentstyle=\color{codegreen},
    keywordstyle=\color{magenta},
    numberstyle=\tiny\color{codegray},
    stringstyle=\color{codepurple},
    basicstyle=\ttfamily\footnotesize,
    breakatwhitespace=false,         
    breaklines=true,                 
    captionpos=b,                    
    keepspaces=true,                 
    numbers=left,                    
    numbersep=5pt,                  
    showspaces=false,                
    showstringspaces=false,
    showtabs=false,                  
    tabsize=2
}
\lstset{style=mystyle}

% --- Configuración de Encabezado y Pie de Página ---
\pagestyle{fancy}
\fancyhf{}
\rhead{Proyecto 2 - DLP}
\lhead{Generador YALex/YAPar}
\rfoot{Página \thepage}

\begin{document}

% ==========================================
% PORTADA
% ==========================================
\begin{titlepage}
    \centering
    \vspace*{1cm}
    
    {\scshape\LARGE Universidad del Valle de Guatemala \par}
    \vspace{1cm}
    {\scshape\Large Diseño de Lenguajes de Programación \par}
    \vspace{1.5cm}
    {\huge\bfseries Proyecto 2: Generador YALex / YAPar \par}
    \vspace{0.5cm}
    {\large\itshape Construcción Autónoma de Analizadores Léxicos y Sintácticos \par}
    \vspace{2cm}
    
    {\Large\bfseries Integrantes del Grupo: \par}
    \vspace{0.5cm}
    {\large 1. Angel Esteban Esquit Hernández - 23221 \par}
    {\large 2. Javier Eduardo España Pacheco - 23361 \par}
    
    \vfill
    {\large Catedrático: Ing. Carlos Jorge Valdez Bautista \par}
    \vspace{0.8cm}
    {\large \today \par}
\end{titlepage}

\newpage
\tableofcontents
\newpage

% ==========================================
% SECCIÓN 1: DISEÑO DEL SOFTWARE
% ==========================================
\section{Diseño del Software}

El sistema \textbf{YALex Studio} es una suite integrada para la creación autónoma de compiladores, la cual consta de dos componentes medulares desarrollados completamente desde cero (sin dependencias externas de librerías de autómatas o motores de expresiones regulares):
\begin{itemize}
    \item \textbf{Generador Léxico (YALex):} Construye analizadores léxicos (lexers) mediante el método directo a partir de archivos de especificación \texttt{.yal}.
    \item \textbf{Generador Sintáctico (YAPar):} Genera analizadores sintácticos (parsers) utilizando la técnica de análisis sintáctico SLR (Simple LR) sobre una colección canónica de items LR(0) a partir de archivos de especificación \texttt{.yapar}.
    \item \textbf{Independencia de los Analizadores Generados:} Los componentes resultantes (\texttt{lexer.py} y \texttt{parser.py}) son completamente autónomos e independientes de los módulos generadores. Esto permite su ejecución autónoma en cualquier entorno de Python estándar sin requerir la presencia de la base de código del compilador.
\end{itemize}

\subsection{Arquitectura General del Sistema}
El proyecto está estructurado de manera modular y desacoplada, separando la lógica de análisis, simulación, generación de código y la interfaz de usuario:

\begin{verbatim}
PRY2-DLP/
+- src/
|  +- yalex_parser/            <- Implementacion original del lexer
|  |  +- parser.py            <- Parser de especificaciones .yal
|  |  +- regex_parser.py      <- Parser de regex y AST
|  |  +- direct.py            <- AFD por metodo directo
|  |  +- dfa.py               <- Minimizacion de AFD
|  |  +- simulator.py         <- Simulacion con traza
|  |  +- codegen.py           <- Generacion del Lexer
|  +- yapar_generator/         <- Generador de parsers LR(0)/SLR
|  |  +- yapar_parser.py      <- Parser de especificaciones .yapar
|  |  +- lr0.py               <- Automata LR(0), closure y transiciones
|  |  +- table.py             <- Generador de tablas SLR
|  |  +- parser.py            <- Maquina de pila de analisis
|  |  +- codegen.py           <- Generacion del Parser
|  +- common/                  <- Utilidades de tokens
|  +- bridge_cli.py            <- Interfaz Bridge
|  +- bridge_yapar.py          <- Integration YAPar
+- desktop-app/                 <- Interfaz de usuario (Tauri + React)
+- examples/                    <- Casos de prueba (Low, Medium, High)
\end{verbatim}

\subsection{Flujo de Trabajo e Integración}
El flujo para la integración del Lexer y el Parser sigue los siguientes pasos:
\begin{enumerate}
    \item \textbf{Entrada YALex (\texttt{.yal}):} Se procesa el archivo de tokens, construyendo el AFD correspondiente para generar un script Python ejecutable (\texttt{lexer.py}).
    \item \textbf{Entrada YAPar (\texttt{.yapar}):} Se procesa la especificación sintáctica, calculando los conjuntos de items LR(0), las transiciones y la tabla SLR para generar un parser autónomo (\texttt{parser.py}).
    \item \textbf{Ejecución en Cadena:} El archivo de código de entrada (\texttt{input.txt}) es tokenizado secuencialmente por el lexer generado, y el flujo de tokens alimenta directamente al parser SLR, que produce el árbol de análisis sintáctico (AST) o reporta los errores sintácticos con precisión.
\end{enumerate}

% ==========================================
% SECCIÓN 2: DOCUMENTACIÓN Y EXPLICACIÓN DEL CÓDIGO
% ==========================================
\section{Documentación y Explicación del Código}

\subsection{Generador de Analizadores Léxicos (\texttt{yalex\_parser})}
Este módulo es el responsable de analizar la especificación de tokens y construir la lógica del escáner:
\begin{itemize}
    \item \textbf{\texttt{parser.py}:} Lee la definición del archivo \texttt{.yal}, separando la sección de definiciones (\texttt{let}), las reglas (\texttt{rule}) y el código empotrado del usuario (header y trailer).
    \item \textbf{\texttt{regex\_parser.py} y \texttt{regex\_ast.py}:} Traducen la notación de las expresiones regulares a un árbol de sintaxis abstracta (AST) autónomo. Soporta operaciones de cerradura de Kleene (\texttt{*}), unión (\texttt{|}), concatenación, y rangos (como \texttt{[a-z]}).
    \item \textbf{\texttt{direct.py}:} Utiliza el método directo para calcular las funciones de árboles de expresiones regulares: \texttt{nullable}, \texttt{firstpos}, \texttt{lastpos} y \texttt{followpos}. Convierte el AST de la regex directamente en un Autómata Finito Determinista (AFD).
    \item \textbf{\texttt{dfa.py}:} Aplica el algoritmo de Hopcroft para minimizar los estados del AFD generado, garantizando eficiencia óptima durante la tokenización.
    \item \textbf{\texttt{simulator.py}:} Contiene el motor que evalúa los caracteres de entrada frente al AFD minimizado utilizando una lógica de \textit{Maximal Munch} (máxima coincidencia).
\end{itemize}

\subsection{Generador de Analizadores Sintácticos (\texttt{yapar\_generator})}
Este módulo implementa la lógica de análisis LALR/SLR basada en la gramática libre de contexto (CFG) provista:
\begin{itemize}
    \item \textbf{\texttt{yapar\_parser.py}:} Analiza el archivo de gramática \texttt{.yapar}. Reconoce tokens declarados con \texttt{\%token}, el símbolo inicial con \texttt{\%start}, y la lista de producciones sintácticas junto con sus acciones semánticas asociadas.
    \item \textbf{\texttt{lr0.py}:} Modela los items LR(0) ($A \rightarrow \alpha \cdot \beta$), estados del autómata, y calcula recursivamente el \texttt{closure} de conjuntos de items y las transiciones mediante la función \texttt{GOTO}. Esto genera el Autómata LR(0) Canónico.
    \item \textbf{\texttt{table.py}:} Computa las funciones auxiliares de gramáticas \texttt{FIRST} y \texttt{FOLLOW} para todos los no-terminales. A partir de los FOLLOW sets y el autómata LR(0), construye la tabla de análisis sintáctico SLR (filas de estados, columnas de terminales y no-terminales) con acciones de:
    \begin{itemize}
        \item \textbf{SHIFT (s):} Desplazar un terminal al tope de la pila y transicionar a un nuevo estado.
        \item \textbf{REDUCE (r):} Reducir una regla de producción aplicando la correspondiente substitución del no-terminal.
        \item \textbf{ACCEPT (acc):} Aceptación exitosa de la cadena de entrada.
        \item \textbf{ERROR:} Estado en blanco en la tabla que representa un error sintáctico.
    \end{itemize}
    \item \textbf{\texttt{parser.py}:} Implementa el motor de ejecución basado en pila (Stack Machine) que simula el autómata SLR, manteniendo el registro exacto del historial de pasos (Trace) para propósitos de depuración.
\end{itemize}

% ==========================================
% SECCIÓN 3: CASOS DE PRUEBA Y CASO COMPLEJO
% ==========================================
\section{Resolución de Casos de Prueba}

El sistema fue validado rigurosamente utilizando tres escenarios de prueba de complejidad incremental ubicados en la carpeta \texttt{examples/}:
\begin{itemize}
    \item \textbf{Calculadora Simple (\texttt{low}):} Valida expresiones aritméticas de sumas y restas utilizando tokens de identificadores y números.
    \item \textbf{Lenguaje Estructurado (\texttt{medium}):} Valida flujos de control simples como condicionales \texttt{if-else}, bucles \texttt{while}, asignaciones y llamadas a funciones.
    \item \textbf{Lenguaje OOP de Alta Complejidad (\texttt{high}):} Representa un subconjunto de lenguaje similar a Java con clases, herencia, constructores, excepciones y arreglos.
\end{itemize}

% ==========================================
% SECCIÓN 4: REPOSITORIO DE GITHUB
% ==========================================
\section{Enlace al Repositorio de GitHub}

El código fuente completo de la implementación, incluyendo el generador de autómatas, el creador de tablas SLR, el parser guiado por pila, las pruebas unitarias y los ejemplos de entrada se encuentra alojado en el siguiente repositorio público de GitHub:

\begin{center}
    \url{https://github.com/Javier-Espana/PRY2-DLP}
\end{center}

\end{document}
